\documentclass[../../main.tex]{subfiles}
\begin{document}

\begin{center}
  {\fontsize{48pt}{58pt}\selectfont \textbf{1978}\normalsize\textbf{年度}}
\end{center}

\begin{table}[htb]
  \centering
  \caption{出題テーマと難易度}
  \begin{tabular}{|c|p{9cm}|c|}
    \hline
    \textbf{問題番号} & \textbf{テーマ} & \textbf{難易度} \\
    \hline
    第1問 & 相加平均と相乗平均の差の不等式 & \\
    \hline
    第2問 & 重心を通る直線と三角形の面積比 & \\
    \hline
    第3問 & 放物線弧と直線に囲まれる図形の二等分 & \\
    \hline
    第4問 & 回転行列の三乗根 & \\
    \hline
    第5問 & 整数根をもつ三次方程式の係数 & \\
    \hline
    第6問 & 三角関数の直交性と期待二乗誤差の最小化 & \\
    \hline
  \end{tabular}
\end{table}

\begin{table}[htb]
  \centering
  \caption{解答}
  \renewcommand{\arraystretch}{1.4}
  \begin{tabular}{|c|c|p{9cm}|}
    \hline
    \textbf{問題番号} & \textbf{小問} & \textbf{答え} \\
    \hline
    第1問 & -- & 不等式は成立し，等号は$c=\sqrt{ab}$のとき \\
    \hline
    第2問 & (1) & $\dfrac{1}{h}+\dfrac{1}{k}=3$ \\
    \cline{2-3}
    & (2) & $\dfrac{4}{9}S\le T\le\dfrac{1}{2}S$ \\
    \hline
    第3問 & (1) & $p=-\dfrac{3}{2}$ \\
    \cline{2-3}
    & (2) & $x=p$は$CD$を二等分する \\
    \cline{2-3}
    & (3) & $x=p$は囲まれた面積を二等分する \\
    \hline
    第4問 & -- &  \\
    \hline
    第5問 & (i) & $(m,n)=(4,5),(0,-3),(-2,-1)$ \\
    \cline{2-3}
    & (ii) & $n=2m-3,-2m-5,-m-3,m+1$のいずれか \\
    \cline{2-3}
    & (iii) & $26$組 \\
    \hline
    第6問 & (1) & $I_{m,n}=\pi$（$m=n$），$0$（$m\ne n$） \\
    \cline{2-3}
    & (2) & 最小は$(m,n)=(2,3),(3,2)$かつ$a=b=\dfrac{5}{16}$のとき \\
    \hline
  \end{tabular}
\end{table}

\paragraph{試験形態}
6問．試験時間は要確認．

\end{document}
